
\begin{center}
	\LARGE
	Uden hjælpemidler
\end{center}
\begin{opgavetekst}{Opgave 1}
\end{opgavetekst}
\begin{delopgave}{1}
	Bestem følgende logaritmer
	\begin{align*}
		&1) \ \log_{5}(125)& &2) \ \log_{7}\left(\frac{1}{49}\right)\\
		&3) \ \log_{2}(1024)& &4) \ \log_{64}(16)\\
		&5) \ \log_{10}(0.00001)& &6) \ \log_{27}(3)\\
		&7) \ \ln(e^{-0.3})& &8) \ \log_{128}(8)\\
	\end{align*}
\end{delopgave}
%
%
%
%
%
%
\begin{opgavetekst}{Opgave 2}
\end{opgavetekst}
\begin{delopgave}{1}
	Bestem den afledede funktion for følgende funktioner
	\begin{align*}
		&1) \ f(x) = 5x^2-7x+4 & &2) \ f(x) = \frac{1}{x} + \sqrt{x}\\
		&3) \ f(x) = 2x^{-0.75}+4\sqrt{x} & &4) \ f(x) = x^{\frac{1}{2}}+4x^{\frac{3}{2}}\\
		&5) \ f(x) = \frac{7}{3}\ln(x) + 0.85 & &6) \ f(x) = 1x + \ln(4) -\sqrt{7}\\
	\end{align*}
\end{delopgave}
%
%
%
%
%
\begin{opgavetekst}{Opgave 3}
	En funktion $f$ er givet ved
	\begin{align*}
		g(x) = 5x^2 - 3x.
	\end{align*}
\end{opgavetekst}
\begin{delopgave}{1}
	Vis ved brug af definitionen for differentialkvotienten, at 
	\begin{align*}
		g'(2) = 17.
	\end{align*}	 
	(Vink: Opskriv først sekanthældningen i $x_0  = 2$ og lad efterfølgende $h \to 0$. )
\end{delopgave}

\newpage
%
%
%
%
%
\begin{opgavetekst}{Opgave 4}
	En funktion $f$ er givet ved
	\begin{align*}
		f(x) = -6x + 2
	\end{align*}
\end{opgavetekst}
\begin{delopgave}{1}
	Vis ved brug af definitionen for differentialkvotienten, at 
	\begin{align*}
		f'(x) = -6 
	\end{align*}
	Vink: Opskriv først sekanthældningen i $x_0$ og lad efterfølgende $h \to 0$. )
\end{delopgave}
%
%
%
%
%
\begin{opgavetekst}{Opgave 5}
\end{opgavetekst}
\begin{delopgave}{1}
	Bestem følgende funktioners ekstremumssteder eller vendetangenter.
	\begin{align*}
		&1) \ f(x) = \ln(x) + x^2 + 5  &   &2) \ f(x) = x^4 - 32x^2 + 10 \\
		&3) \ f(x) = \frac{1}{3}x^3 + \frac{1}{2}x^2 - 6x + 4  &   &4) \ f(x) = \frac{2}{5}x^5-27x^2
	\end{align*}
\end{delopgave}
%
%
%
%
%
\begin{opgavetekst}{Opgave 6}
	En funktion $f$ er givet ved
	\begin{align*}
		f(x) = x^3 - x^2 + 3x - 15
	\end{align*}
\end{opgavetekst}
\begin{delopgave}{1}
	Bestem hældningen af tangenten til $f$ i $x = 1$ og bestem en ligning for tangenten i punktet $(1,f(1))$
\end{delopgave}
\begin{meretekst}
	Det oplyses, at der er ét andet punkt $(x_0,f(x_0))$, hvor tangenten er parallel til tangenten i punktet $(1,f(1))$.
\end{meretekst}
\begin{delopgave}{2}
	Bestem tallet $x_0$.
\end{delopgave}
%
%
%
%
%
\begin{opgavetekst}{Opgave 7}
	Et polynomium $p$ er givet ved
	\begin{align*}
		p(x) = ax^2 + bx + c
	\end{align*}
\end{opgavetekst}
\begin{delopgave}{1}
	Vis, at $p$ har toppunkt i $\left(\frac{-b}{2a}, \frac{-d}{4a}\right)$, hvor $d$ er diskriminanten for $p$
\end{delopgave}